% Extracto íntegro por residencia del propietario científico definitivo de masas.
\section{Secteur nul et positivité}

La nullité ne s'obtient pas en prenant une limite du caractère positif. Le domaine
se sépare d'abord en
\[
 \mathcal H_{\rm phys}=\mathcal H_0\oplus\mathcal H_+,
\]
avec
\[
 \mathcal M|_{\mathcal H_0}=0,
 \qquad
 \mathcal M|_{\mathcal H_+}>0.
\]
Le photon appartient à la carte nulle électromagnétique et les gluons à la carte
nulle adjointe de couleur. Leurs différentes représentations, polarisations et
contraintes de jauge restent visibles bien que les deux masses soient nulles.

La positivité exclut une espèce élémentaire avec \(m^2<0\) dans le
codomaine massique. Une valeur propre négative du hessien d'un fond ne définit pas un
tachyon comme particule : elle diagnostique l'instabilité du fond et exige de changer
la solution autour de laquelle on linéarise.

\section{Couleur, saveurs de quarks et mélange CKM}

L'orbite chromatique se réalise dans le qutrit de couleur. Les opérateurs de Weyl
engendrent \(\mathfrak{sl}_3(\mathbb C)\) ; la sélection antihermitienne de trace
nulle produit \(\mathfrak{su}(3)\), et la connexion finie transporte les trois
cartes chromatiques. La masse de saveur doit être invariante sous cette action.

Soit \(U_{\rm col}:SU(3)\to\mathcal U(V_{\rm col})\). La réduction chromatique
est l'espérance de Reynolds
\[
 \mathbb E_{\rm col}(T)
 =\int_{SU(3)}U_{\rm col}(g)T U_{\rm col}(g)^*\,dg.
\]
Dans la représentation fondamentale irréductible,
\[
 \mathbb E_{\rm col}(\mathcal M_q)
 =m_q(\mu,\mathscr R)I_3.
\]
La masse publiée est donc indépendante du rouge, du vert et du bleu. Les
coordonnées \(\mu\) et \(\mathscr R\) identifient l'échelle et le schéma de
renormalisation ; ce ne sont pas des ambiguïtés cachées et elles doivent toujours accompagner le
nombre.

La dépendance en échelle est une section d'un fibré sur
\(t=\log\mu\). Si \(\gamma_m(g(t))\) est l'endomorphisme induit par la
connexion chromatique, le transport est défini par
\[
 \nabla_{\rm RG}m_q
 =\left(\frac{d}{dt}+\gamma_m(g(t))\right)m_q=0,
\]
\[
 m_q(t_2)
 =\mathcal P\exp\!\left(-\int_{t_1}^{t_2}
 \gamma_m(g(t))\,dt\right)m_q(t_1).
\]
Deux valeurs appartenant à des échelles ou à des schémas distincts ne sont pas des nombres
comparables tant qu'elles ne sont pas transportées vers la même fibre. Cette équation clôt le
type mathématique de la masse courante : ce n'est pas un scalaire absolu isolé, mais
une section transportée de manière covariante.

Pour les mésons et les baryons, les projecteurs de singulet sont
\[
 P_{q\bar q}^{(1)}
 =\frac1{3}|\Omega\rangle\langle\Omega|,
 \qquad
 |\Omega\rangle=\sum_{a=1}^{3}|a\bar a\rangle,
\]
et
\[
 P_{qqq}^{(1)}
 =|\varepsilon\rangle\langle\varepsilon|,
 \qquad
 |\varepsilon\rangle
 =\frac1{\sqrt6}\sum_{a,b,c}\varepsilon_{abc}|abc\rangle.
\]
Les deux éliminent l'étiquette chromatique observable sans effacer l'orbite interne
qui est intervenue dans la construction.

Le mélange CKM appartient à une application distincte. Si \(B_u\) et \(B_d\) sont les
opérateurs de masse dans les bases de type haut et de type bas, la matrice unitaire
interne satisfait
\[
 B_d^{(u)}=V_{\rm CKM} B_d V_{\rm CKM}^*.
\]
Les angles dérivés de \(A\), \(C^*\) et \(\Delta_4\), la phase
\(\delta_{\rm CP}\), l'invariant de Jarlskog et l'identité du déterminant
du commutateur déterminent le mélange sans créer les valeurs propres de masse. Ainsi
restent séparés le spectre, la couleur, la génération et le changement de base.

\section{Neutrinos, mélange PMNS et critère de Majorana}

Les neutrinos sont des sections fermioniques électriquement neutres. Le
caractère fermionique descend de l'holonomie spinorielle, de la complétion extérieure et
des relations canoniques d'anticommutation ; la neutralité n'altère pas leur
statistique. Le support interne est
\[
 \mathcal H_\nu=
 \bigoplus_{j=1}^{4}\mathcal H_{\nu,G_j},
 \qquad
 \dim\mathcal H_\nu=2+4+6+8=20.
\]
Les quatre profondeurs structurelles ne sont ni des saveurs ni des valeurs propres. Le
sous-espace physique tridimensionnel s'obtient au moyen du projecteur neutre
\(P_\nu\), et l'opérateur de masse est
\[
 M_\nu=P_\nu\mathcal M_\nu P_\nu.
\]
Ses trois valeurs propres ordonnées sont les masses ; l'échelle absolue provient de la
même horloge massique qui fixe les autres branches, et non des angles de mélange.

Soit \(M_\ell=P_\ell\mathcal M_\ell P_\ell\) l'opérateur chargé. Notons
\(F_\ell\) et \(F_\nu\) des repères spectraux orthonormaux complets, ordonnés par
les projecteurs spectraux de \(M_\ell\) et \(M_\nu\). La matrice de mélange est
\[
 U_{\rm PMNS}=F_\ell^*F_\nu.
\]

\begin{theorem}[Clôture spectrale du mélange leptonique]
Si les deux repères sont complets, \(U_{\rm PMNS}\) est unitaire. Si les
spectres sont simples, elle est déterminée à des réajustements diagonaux de phase et à des
permutations fixées par l'ordre spectral près. S'il existe une dégénérescence, le
résultat exact est une classe double
\[
 \prod_\lambda U(m_{\ell,\lambda})
 \backslash U(3) /
 \prod_\mu U(m_{\nu,\mu}),
\]
et non une matrice numérique artificiellement sélectionnée.
\end{theorem}

\begin{proof}
L'orthonormalité donne
\(U_{\rm PMNS}^*U_{\rm PMNS}=F_\nu^*F_\ell F_\ell^*F_\nu=I_3\).
Les changements de base à l'intérieur de chaque sous-espace dégénéré agissent par la
gauche ou par la droite, ce qui produit exactement la classe double
indiquée.
\end{proof}

Le registre complet détermine les éléments de \(M_\nu\) au moyen de la retenue,
du feuillet, de l'orientation, de la lacune et de l'holonomie. De manière équivalente, on peut employer
le noyau hermitien
\[
 K_{\rm reg}(c,d)
 =\sum_{r\in\{\rm car,hoj,ori,vac,hol\}}
 \frac{\langle\Xi_r(c),\Xi_r(d)\rangle}
      {\sum_x\|\Xi_r(x)\|^2},
\]
en omettant tout terme de norme nulle. Sa matrice de Gram est positive et
construit l'opérateur avant la diagonalisation. La phase de rang un
\[
 I+(e^{i\pi/6}-1)|b_K\rangle\langle b_K|
\]
demeure un contrôle négatif : son unitarité formelle ne suffit pas et son angle
\(\theta_{13}\) réfute son identification au mélange physique de rang
complet.

Soit \(G_0\) la matrice tridimensionnelle obtenue en agrégeant le noyau précédent
sur les trois secteurs neutriniques. La condition finie
\[
 \det G_0\ne0
\]
est le certificat exact que les observables incluses dans le registre
séparent trois directions. Dans ce cas, le facteur polaire
\[
 U_{G_0}=G_0(G_0^*G_0)^{-1/2}
\]
est unitaire et de rang \(3\). Si \(\det G_0=0\), le noyau déclaré ne
détermine pas de transport tridimensionnel : il faut ajouter une observable interne
indépendante et recalculer la matrice de Gram. Il n'est pas licite d'imposer
l'inversibilité en ajoutant un terme construit à partir de la matrice de mélange
que l'on cherche à obtenir. La matrice physique est fixée par les repères
spectraux de \(M_\ell\) et \(M_\nu\) ; le critère de Gram certifie la suffisance
du registre, sans sélectionner à lui seul ces repères.

Soit \(\mathcal C_\nu\) l'autoconjugaison antiunitaire du secteur neutre. Le
critère de Majorana est
\[
 \mathcal C_\nu^2=I,
 \qquad
 \mathcal C_\nu M_\nu=M_\nu\mathcal C_\nu,
 \qquad
 \mathcal C_\nu H_\nu=H_\nu\mathcal C_\nu,
\]
avec l'existence d'une base réelle fixe dans la réalification du
sous-espace physique. Lorsque ces conditions sont remplies, chaque mode peut
être représenté par \(\psi=\mathcal C_\nu\psi\). La branche de Dirac exige, en
revanche, une charge conservée \(Q\) telle que
\[
 \mathcal C_\nu Q\mathcal C_\nu^{-1}=-Q,
 \qquad Q\psi\ne0,
\]
et conserve comme modes distincts \(\psi\) et \(\mathcal C_\nu\psi\). Si aucun
des deux ensembles de conditions n'est établi, le type reste indécidé.
La neutralité, l'involution de mot ou le secteur d'Ising, pris
séparément, ne tranchent pas cette alternative.

Une profondeur structurelle supplémentaire ne constitue pas à elle seule un neutrino
stérile. Un secteur stérile candidat est défini par un projecteur non nul
\(P_s\) qui satisfait
\[
 P_s P_{\rm act}=0,
 \qquad P_s J_{\text{faible}}=0,
\]
avec une dynamique définie sur \(P_s\mathcal H_\nu\). Si, de plus,
\([P_s,M_\nu]=0\), le secteur est réducteur et ne se mélange pas au secteur actif. Si
\(P_sM_\nu P_{\rm act}\ne0\), il existe un mélange actif--stérile et nécessairement
\([P_s,M_\nu]\ne0\). Par conséquent, l'existence du secteur et l'existence du
mélange sont des propriétés distinctes. Si le projecteur faiblement neutre n'existe pas,
le canal grossier supplémentaire est une profondeur interne du support et non une quatrième
espèce.

\section{Bosons massifs et secteur scalaire}

Les secteurs \(W\), \(Z\) et scalaire possèdent leurs propres descripteurs de
représentation, de route et de signature. Leur réalisation n'exige pas de les faire entrer de force dans une
fibre élémentaire de rang un. On construit d'abord l'opérateur positif de
fibre puis, dans un sous-espace \(P_X\mathcal H_X\) de dimension finie,
l'opérateur effectif de canaux
\[
 H_{X,\rm eff}(z)
 =P_X H_X P_X
 +P_X H_X Q_X(z-Q_X H_X Q_X)^{-1}Q_X H_X P_X.
\]
Si la résolvante réduite admet un prolongement méromorphe à travers la coupure du
continu, les pôles sont les zéros isolés, comptés avec multiplicité, de
\[
 \det\bigl(z-H_{X,\rm eff}(z)\bigr)=0,
 \qquad
 z_X=M_X c_{\rm int}^2-\frac{i}{2}\Gamma_X.
\]
Un canal ouvert ne garantit pas à lui seul l'existence de tels zéros. Lorsqu'il
existe un zéro avec le résidu physique prescrit, la partie réelle publie la masse de
pôle et la partie imaginaire la largeur. Pour le
secteur scalaire, l'identification d'une excitation exige en outre le projecteur
scalaire et ses couplages de jauge et fermioniques. Deux pôles indépendants
déterminent deux états ; une seconde étiquette sans projecteur ni pôle ne crée pas une
particule.

\section{États composés et loi de liaison}

Pour des constituants \(X_1,\ldots,X_n\), l'entrée est le registre composé
\[
 \mathfrak C_{12}^{(n)}:
 (\Gamma_1,\ldots,\Gamma_n;\mathcal T,\partial)
 \longmapsto\mathcal L_{\rm comp},
\]
et non la somme des masses isolées. La composition binaire
\[
 L_1\star L_2=L_1+L_2+\mathfrak b(L_1,L_2)
\]
est associative exactement lorsque
\[
 \mathfrak b(L_1,L_2)+\mathfrak b(L_1\star L_2,L_3)
 =\mathfrak b(L_2,L_3)+\mathfrak b(L_1,L_2\star L_3).
\]
La signature, le raffinement et l'opérateur composé sont calculés après la
liaison :
\[
 \mathcal L_{\rm comp}\longmapsto\sigma_{\rm comp}
 \longmapsto\eta_{\rm comp}\longmapsto\mathcal M_{\rm comp}.
\]
Le projecteur physique réunit le singulet de couleur, le spin et la parité, l'isospin, la charge,
la symétrie d'échange et le canal. La masse stable est une valeur propre de la
compression ; lorsque les hypothèses de prolongement méromorphe sont remplies, la masse
et la largeur d'une résonance proviennent d'un pôle de Feshbach. Ainsi,
les mésons, les baryons, les états exotiques, les noyaux, les atomes et
les molécules possèdent le même schéma causal sans se réduire à une somme de masses.

Deux arbres de constituants représentent la même espèce s'il existe un
unitaire qui entrelace simultanément leurs opérateurs, leurs projecteurs de canal et leurs
pôles. En l'absence d'un tel unitaire, ils sont conservés comme des réalisations distinctes,
bien qu'ils partagent des constituants et des nombres quantiques partiels.

\section{Résonances, largeurs et survie}

Soit \(S_X\) le projecteur sur le sous-espace survivant et
\(T_X=S_X U S_X\) l'opérateur comprimé d'un pas. S'il possède une valeur propre
dominante simple
\[
 \lambda_X=r_X e^{-i\theta_X},\qquad 0<r_X<1,
\]
fixons la branche \(\theta_X\in I_\theta\), où \(I_\theta\) est l'intervalle
fondamental de quasi-énergie déterminé par la carte temporelle et son orientation.
Alors
\[
 M_X=\frac{\hbar_{\rm int}}{t_0c_{\rm int}^2}\theta_X,
 \qquad
 \Gamma_X=-\frac{2\hbar_{\rm int}}{t_0}\log r_X,
 \qquad
 \tau_X=\frac{\hbar_{\rm int}}{\Gamma_X}.
\]
Sans le choix de \(I_\theta\), la phase ne détermine \(\theta_X\) que modulo
\(2\pi\) et la masse de quasi-énergie n'est pas un scalaire bien défini.
La largeur n'est pas une sixième coordonnée de la signature massique : elle provient du module
du mode ouvert. L'égalité de masse et de largeur entre particule et
antiparticule exige de transporter également les canaux et la prescription causale, et pas seulement
la parité de la signature.

\section{Complétion multiparticulaire et statistique}

Soit \(\mathsf{Hol}_{\pm}^{\leq1}\) la catégorie monoïdale de fibres à une
particule, avec des morphismes de raffinement contractifs, en particulier unitaires,
et une parité d'holonomie
\(\epsilon\in\{+1,-1\}\). La complétion de Fock est définie par le foncteur
\[
 \mathfrak F(H,\epsilon)=
 \begin{cases}
 \displaystyle\bigoplus_{n\ge0}\operatorname{Sym}^n(H),
     &\epsilon=+1,\\[2mm]
 \displaystyle\bigoplus_{n\ge0}\bigwedge^n\!H,
     &\epsilon=-1,
 \end{cases}
\]
et, pour un entrelaceur contractif \(T:H\to H'\),
\[
 \mathfrak F(T)=
 \bigoplus_{n\ge0}\operatorname{Sym}^n(T)
 \quad\text{ou}\quad
 \bigoplus_{n\ge0}\bigwedge^n\!T.
\]

\begin{theorem}[Transport interne de la statistique]
La complétion précédente est un foncteur d'opérateurs bornés sur
\(\mathsf{Hol}_{\pm}^{\leq1}\), compatible avec le raffinement et
monoïdal gradué. Dans la branche \(\epsilon=-1\), les opérateurs de création et
d'annihilation satisfont CAR et les opérateurs d'occupation vérifient
\(N_s^2=N_s\) ; dans la branche \(\epsilon=+1\), ils satisfont CCR. Par conséquent,
l'exclusion de Pauli et les statistiques fermionique et bosonique découlent de la
parité d'holonomie une fois fixée la représentation à une particule.
\end{theorem}

\begin{proof}
Les puissances symétriques et extérieures préservent les compositions et les identités ;
la contractivité de \(T\) garantit que leur somme directe définit un opérateur
borné. Pour les morphismes non contractifs, la seconde quantification exige un
domaine dense explicite et n'appartient pas à cette catégorie d'opérateurs bornés.
La règle des signes du produit extérieur produit CAR et l'idempotence de
l'occupation ; le quotient symétrique produit CCR. La seconde quantification des
morphismes de raffinement entrelace les opérateurs correspondants.
\end{proof}

Cette démonstration clôt le transport algébrique interne. L'identification
au théorème relativiste de spin--statistique exige, en outre, le pont
lorentzien et la localité déclarés en mécanique dimensionnelle ; les deux énoncés
ne doivent pas être confondus.

Pour un Hamiltonien comprimé \(H_{X,V}\) sur un volume ou un raffinement
fini \(V\), l’état thermique est complètement défini par
\[
 Z_{X,V}(\beta)=\operatorname{Tr}e^{-\beta H_{X,V}},
 \qquad
 \rho_{X,V,\beta}=Z_{X,V}(\beta)^{-1}e^{-\beta H_{X,V}}.
\]
La limite thermodynamique existe lorsque la pression
\[
 p_X(\beta)=\lim_{V\nearrow\infty}
 \frac1{\beta |V|}\log Z_{X,V}(\beta)
\]
existe et est indépendante de la suite cofinale. La condensation bosonique se
décide par la saturation de la densité d'états excités ; elle n'est pas proclamée à
partir de la seule existence d'une complétion symétrique. Ainsi, le problème
est exprimé au moyen d'une condition mathématique vérifiable pour le
Hamiltonien et la limite, au lieu d'une attribution nominale.

\section{Secteurs topologiques et états virtuels}

Dans un secteur topologique, le codomaine primaire est une catégorie tressée, et non
une droite de masses. Pour des objets simples \(a,b,c\), les données physiques comprennent
\[
 N_{ab}^{c},\qquad F^{abc}_{d},\qquad R^{ab}_{c},
\]
avec des équations pentagonale et hexagonale. Une masse scalaire n'existe que
lorsqu'une réalisation hamiltonienne fournit un projecteur spectral et un
gap. Les secteurs Fibonacci, Ising/Majorana et \(\mathbb Z_6\) sont ainsi
clos comme codomaines catégoriques ; aucune valeur numérique universelle ne leur est attribuée
par analogie avec une particule élémentaire.

Dans le secteur \(\mathbb Z_6\), la fusion et le tressage sont
\[
 a\otimes b=a+b\pmod 6,
 \qquad
 R_{ab}=\zeta_6^{ab},
 \qquad \zeta_6=e^{2\pi i/6}.
\]
Pour Fibonacci,
\[
 \tau\otimes\tau=\mathbf 1\oplus\tau,
 \qquad
 F_\tau^{\tau\tau\tau}=
 \begin{pmatrix}
 \varphi^{-1}&\varphi^{-1/2}\\
 \varphi^{-1/2}&-\varphi^{-1}
 \end{pmatrix},
\]
\[
 R_{\mathbf1}^{\tau\tau}=e^{-4\pi i/5},
 \qquad
 R_\tau^{\tau\tau}=e^{3\pi i/5}.
\]
Le secteur d'Ising possède des objets \(\mathbf1,\psi,\sigma\) et les règles
\[
 \psi\otimes\psi=\mathbf1,
 \qquad
 \psi\otimes\sigma=\sigma,
 \qquad
 \sigma\otimes\sigma=\mathbf1\oplus\psi,
\]
\[
 F_\sigma^{\sigma\sigma\sigma}
 =\frac1{\sqrt2}
 \begin{pmatrix}1&1\\1&-1\end{pmatrix},
 \qquad
 R_{\mathbf1}^{\sigma\sigma}=e^{-\pi i/8},
 \qquad
 R_\psi^{\sigma\sigma}=e^{3\pi i/8}.
\]
Ces matrices satisfont les relations de cohérence dans les conventions
indiquées. Le dédoublement
\(\mathcal C\boxtimes\mathcal C^{\rm rev}\) conserve les deux orientations et
évite d'introduire une chiralité par un choix externe.

Un état virtuel est une section de persistance finie dans une amplitude ou une
résolvante. Il n'appartient pas au spectre asymptotique et n'exige pas de masse de pôle
réelle. Ses paramètres sont la virtualité, le canal et la durée de
persistance. Cette définition évite de transformer une ligne interne de calcul en
une nouvelle espèce.

\Needspace{25\baselineskip}
\section{Gravitation, torsion et mode de spin 2}

La gravitation entre dans la mécanique dimensionnelle par la connexion, la courbure et la
torsion. Elle ne requiert pas de particule élémentaire de spin \(2\) comme donnée de
départ. Soit \(\mathcal S_{\rm grav}\) l'action géométrique et
\(\mathcal H_{\rm TT}\) le sous-espace transverse et sans trace de son Hessien
en une solution. L'observable linéarisée est
\[
 \mathcal K_{\rm TT}
 =P_{\rm TT}\,\delta^2\mathcal S_{\rm grav}\,P_{\rm TT}.
\]
Un mode collectif de spin \(2\) n'existe comme particule asymptotique que
si la résolvante de \(\mathcal K_{\rm TT}\) possède un pôle physique de résidu
positif et respecte les contraintes. Si un tel pôle n'existe pas, l'interaction
gravitationnelle demeure intégralement dans la connexion et la torsion. Ainsi, le
graviton n'est pas un élément nécessaire de la loi des masses, et son apparition éventuelle
est décidée par un critère spectral, et non par une case nominale.
